%=============================================================================
% WAVE 3, ITERATION 8: DEFINITIVE PATENT TRIAGE AND TECHNOLOGY ROADMAP
%
% Agent: Technology & Patent Assessment (W3i8)
% Date: 20 March 2026
%
% W3i7 ESTABLISHED (cumulative through 7 iterations):
%   P1: DEAD (zero resources)
%   P2: Passive detection only; active invalidated
%   P3: QEC advantages survive (7.7-83x); 250 dB gravitational immunity
%   P4: Underground WPT robust but LCOE ~27c/kWh (conversion bottleneck)
%   P5: 11/13 predictions active; Hubble tension robust
%   P6: Gap-filler not competitor; SQL 9.1um adequate
%   P7: Ultra-fast pulsed buffer only (DEW, fusion, accelerator)
%   P8: EFFECTIVELY DEAD (planetary mass required)
%   P9: All claims confirmed at SQL
%   P10: BQP-complete; 10^12x simulation speedup
%   P11: Detection roadmap defined (3 phases)
%   P12: 20% efficiency, ~27c/kWh; niche markets only
%   P13: CLONETS-DS 1.0 day; GPS diurnal archival test
%   P14: Canonical 2-component framework; 1 free parameter
%   P15: Channel B highest priority
%
% W3i8 MANDATE:
%   1. Definitive one-paragraph triage for all 15 patents
%   2. Complete technology roadmap: what gets built, in what order,
%      at what cost, and what each experiment tests
%=============================================================================

\documentclass[11pt]{article}
\usepackage{amsmath,amssymb,amsthm}
\usepackage{geometry}
\geometry{margin=1in}
\usepackage{booktabs}
\usepackage{enumitem}
\usepackage{hyperref}
\usepackage{xcolor}
\usepackage{tcolorbox}
\usepackage{multirow}
\usepackage{array}
\usepackage{siunitx}
\usepackage{float}
\usepackage{longtable}
\usepackage{pdflscape}

\newtheorem{theorem}{Theorem}
\newtheorem{proposition}{Proposition}
\newtheorem{finding}{Finding}
\newtheorem{correction}{Correction}
\newtheorem{result}{Result}

\newcommand{\ee}[1]{\times 10^{#1}}
\newcommand{\psigrav}{\psi_{\rm grav}}
\newcommand{\dpsiEM}{\delta\psi_{\rm EM}}
\newcommand{\lambdaone}{|\lambda{-}1|}
\newcommand{\Hzero}{H_0}
\newcommand{\aM}{a_0}
\newcommand{\GN}{G_N}
\newcommand{\Qpsi}{Q_\psi}
\newcommand{\Meff}{M_{\rm eff}}

\tcbuselibrary{skins,breakable}

\title{\textbf{Wave 3 Iteration 8: Definitive Patent Triage}\\[6pt]
\textbf{and Complete Technology Roadmap}\\[12pt]
{\large All 15 DFD Patents: Status, Value, Resources, Dependencies}\\[4pt]
{\large Experimental Programme: Phase~0 through Phase~IV}}
\author{DFD Research Programme --- Technology \& Patent Assessment Agent (W3i8)}
\date{20 March 2026}

\begin{document}
\maketitle
\tableofcontents
\newpage

%=============================================================================
\section*{Executive Summary}
%=============================================================================

\begin{tcolorbox}[colback=blue!5!white, colframe=blue!60!black,
  title=\textbf{W3i8 PATENT TRIAGE: READ THIS FIRST}]

After seven iterations of adversarial scrutiny across 17 parallel agents,
the 15 DFD patents sort into four tiers:

\medskip
\textbf{TIER 1 --- ALIVE (active development warranted):}
\begin{itemize}[nosep]
  \item P5 (Cosmological Predictions): 11/13 testable; Hubble tension robust
  \item P9 (Psi-Field Positioning): SQL-sufficient; transformative underwater
  \item P11 (EM Coupling / SRF Detection): Gateway experiment for all of DFD
  \item P14 (Canonical Framework): $\alpha^{57}$ confirmed; 1 free parameter
\end{itemize}

\textbf{TIER 2 --- NICHE (viable in restricted domains):}
\begin{itemize}[nosep]
  \item P2 (Resonators/Metamaterials): Passive detection only; LIGO 6th channel
  \item P3 (Quantum Coherence): 7.7--83$\times$ qubit advantage confirmed
  \item P4 (Underground WPT): Robust physics; blocked by conversion bottleneck
  \item P6 (Decihertz GW): Gap-filler; ground-based decihertz access
  \item P7 (Psi-Storage): Ultra-fast pulsed buffer for DEW/fusion/accelerators
  \item P10 (Field Computing): $10^{12}\times$ simulation speedup (BQP-standard)
  \item P13 (Clock Comparison): CLONETS-DS 1-day detection; archival tests
  \item P15 (Generator): Channel~B highest experimental priority
\end{itemize}

\textbf{TIER 3 --- THEORETICAL ONLY:}
\begin{itemize}[nosep]
  \item P12 (City-Scale WPT): Economics killed at 20\% efficiency
\end{itemize}

\textbf{TIER 4 --- DEAD:}
\begin{itemize}[nosep]
  \item P1 (Field Drag / Cloaking): Signal $10^{15}\times$ below detection
  \item P8 (Psi-Communications): Mass gap $>10^{11}\times$; no workaround
\end{itemize}

\medskip
\textbf{The single most important experiment}: SQMS Phase~I (P11),
\$2--4M, 2027--2028. Everything downstream depends on detecting a
$\psi$-field signature.

\textbf{The single most important theoretical milestone}: Breaking the
50\% SRF conversion bottleneck. This gates P4, P7, and P12 economics.

\end{tcolorbox}

%=============================================================================
%=============================================================================
\part{Definitive Patent Triage: All 15 Patents}
\label{part:triage}
%=============================================================================
%=============================================================================

%=============================================================================
\section{Patent 1: Field Drag / Gravitational Cloaking}
\label{sec:P1}
%=============================================================================

\begin{tcolorbox}[colback=red!5!white, colframe=red!60!black,
  title=\textbf{P1: DEAD --- Resource Allocation: 0/5}]
\textbf{Status: DEAD.}
The field drag signal from laboratory-mass objects is $\sim 10^{-27}$~s
timing shift, which is $10^{15}\times$ below the detection threshold of any
foreseeable clock technology. The two-component framework (W3i5) resolved
the earlier T5 self-consistency issue, confirming P1 is theoretically
self-consistent but experimentally inaccessible. All active claims
(gravitational shielding, field drag propulsion, cloaking) require
$\lambdaone$ values many orders of magnitude above the current bound
($\lambdaone < 1.56\ee{-9}$). No surviving claim has experimental
relevance. The most valuable aspect of P1 is its role in establishing
the DFD metric structure ($ds^2 = -e^{2\psi}c^2dt^2 + e^{-2\psi}dx^2$),
which is now part of the P14 canonical framework.

\textbf{Most valuable claim:} Theoretical --- the DFD metric coupling
to matter (inherited by P14).

\textbf{Dependencies:} None active. P1 depends on P14 (framework) but
contributes nothing to other patents.
\end{tcolorbox}

%=============================================================================
\section{Patent 2: Resonators and Metamaterials}
\label{sec:P2}
%=============================================================================

\begin{tcolorbox}[colback=yellow!5!white, colframe=yellow!60!black,
  title=\textbf{P2: NICHE --- Resource Allocation: 2/5}]
\textbf{Status: NICHE.}
Active resonator/metamaterial claims are invalidated: the $\lambdaone$
coupling is too weak ($< 1.56\ee{-9}$) for any material to produce
measurable active effects. However, P2 survives in two critical roles:
(1) \emph{passive detection} --- P2 resonator technology underpins the
LIGO 6th-channel search and the SRF cavity programme (P11), where Nb
passive FOM dominates; (2) \emph{P15 Channel~B} --- the toroidal cavity
geometry from P2 is essential for the highest-priority experimental
channel. The W3i7 LIGO correction ($\lambdaone^{\rm LIGO}_{\rm corrected}
\sim 3.5\ee{-14}$, not $6\ee{-19}$) weakens the LIGO reach but the
analysis is still valuable. The Nb cavity preparation protocol
(medium-$T$ N-doping, EP, low-$T$ bake) is directly applicable to P11
Phase~I.

\textbf{Most valuable claim:} Passive Nb SRF cavity as $\psi$-field
detector (enables P11 and P15).

\textbf{Dependencies:} P2 $\to$ P11 (cavity technology); P2 $\to$ P15
(toroidal geometry for Channel~B). P2 depends on P14 (coupling physics).
\end{tcolorbox}

%=============================================================================
\section{Patent 3: Quantum Coherence Control (Psi-QEC)}
\label{sec:P3}
%=============================================================================

\begin{tcolorbox}[colback=green!5!white, colframe=green!40!black,
  title=\textbf{P3: NICHE --- Resource Allocation: 2/5}]
\textbf{Status: NICHE (high theoretical value, contingent on detection).}
The psi-QEC code family $[[2K{+}1,1,K{+}1]]$ survives all W3i7 scrutiny.
Transversal gates are Clifford-only (Eastin--Knill), requiring magic
state distillation for the $T$ gate, but the overhead is modest (105
physical qubits per $T$ gate vs.\ 810 for surface codes). The net
qubit advantage is 7.7$\times$ ($T$-heavy) to 83$\times$
(Clifford-heavy), confirming the 10--28$\times$ range for typical
algorithms. The 250~dB gravitational noise immunity and $67^K$
error suppression are unique. The circuit-level noise model is fully
specified with threshold $p_{\rm th} \geq 1.1\ee{-2}$. However, all
P3 advantages are \emph{contingent on the psi-field existing}: until
P11 detects the field, P3 is theoretical. No immediate resource
allocation is needed beyond monitoring quantum error correction
literature.

\textbf{Most valuable claim:} $67^K$ error suppression factor from
psi-gravitational noise immunity (unique to DFD).

\textbf{Dependencies:} P3 depends on P11 (detection), P14 (field
existence). P3 $\to$ P10 (enables quantum computing advantage).
\end{tcolorbox}

%=============================================================================
\section{Patent 4: Underground Wireless Power Transfer}
\label{sec:P4}
%=============================================================================

\begin{tcolorbox}[colback=yellow!5!white, colframe=yellow!60!black,
  title=\textbf{P4: NICHE --- Resource Allocation: 2/5}]
\textbf{Status: NICHE (physics robust; economics blocked by conversion).}
The underground WPT corridor concept is physically sound: $\dpsiEM$
propagates through arbitrary matter with 81\% transport efficiency over
93.5~km ($d^* = 93.5$~km at $Q_\psi = 10^9$). W3i7 confirmed that
parametric pumping via $\nabla\psigrav$ is negligible at terrestrial
gravity ($\gamma_{\rm par} = 3.35\ee{-14}$~m$^{-1}$, gain $\sim 10^{-13}$
even with resonant enhancement). The critical correction from W3i7:
the W3i6 LCOE of \$0.004/kWh \emph{omitted the two conversion stages}
(EM$\to\dpsiEM$ at 50\%, $\dpsiEM\to$EM at 50\%). Corrected end-to-end
efficiency is 18.8\%, giving LCOE $\sim$27~c/kWh at \$0.05/kWh input.
P4 and P12 share the same conversion bottleneck. Viable only for
niche markets where conventional delivery costs exceed $\sim$25~c/kWh
(deep undersea, deep mining, space).

\textbf{Most valuable claim:} Matter-penetrating wireless energy
corridor --- no physical conduit required.

\textbf{Dependencies:} P4 depends on P11 (detection), P15/P2 (SRF
conversion technology). P4 $\to$ P12 (network topology uses P4
corridors). Breaking the 50\% conversion bottleneck is the gate.
\end{tcolorbox}

%=============================================================================
\section{Patent 5: Cosmological Predictions}
\label{sec:P5}
%=============================================================================

\begin{tcolorbox}[colback=green!5!white, colframe=green!40!black,
  title=\textbf{P5: ALIVE --- Resource Allocation: 4/5}]
\textbf{Status: ALIVE (highest scientific value).}
11 of 13 predictions remain active after W3i7. The Hubble tension
prediction ($\Delta H_0 = 3.9 \pm 1.1$~km/s/Mpc via void outflow)
is robust under two-component: no quantity depends on $\dpsiEM$.
The Pioneer anomaly prediction was demoted from Tier~2 to Tier~3
(ungrounded coincidence: no first-principles mechanism produces
exactly $c H_0$), but this does not affect the 11/13 count since
Pioneer was never high-priority. Five predictions are testable from
archival data \emph{right now}: GPS diurnal drift, Cassini ranging
residuals, pulsar timing anomalies, galaxy rotation curve shape
parameters, and the Tully--Fisher zero-point. T4 (Cold Spot
over-prediction) was reduced from $3$--$5\times$ to $1$--$2\times$
by the void-evolution cancellation effect discovered in W3i7.
P5 is the primary \emph{scientific} case for DFD.

\textbf{Most valuable claim:} Simultaneous resolution of the Hubble
tension, MOND phenomenology, and dark energy from a single scalar field.

\textbf{Dependencies:} P5 depends on P14 (canonical framework), P13
(clock tests verify sector~A). P5 is independent of all technology
patents. P5 $\to$ all (scientific credibility drives funding).
\end{tcolorbox}

%=============================================================================
\section{Patent 6: Decihertz Gravitational Wave Detection}
\label{sec:P6}
%=============================================================================

\begin{tcolorbox}[colback=yellow!5!white, colframe=yellow!60!black,
  title=\textbf{P6: NICHE --- Resource Allocation: 1/5}]
\textbf{Status: NICHE (gap-filler, not competitor).}
W3i7 established definitively that DFD P6 does \emph{not} beat LISA,
DECIGO, BBO, ET, or CE on raw strain sensitivity at \emph{any}
frequency. The advantage is contextual: no funded detector covers
0.01--1~Hz (the ``decihertz gap''). At $N \geq 10^4$ sensors, P6
provides ground-based decihertz access using demonstrated atom
interferometry technology (MAGIS/AION). The only astrophysical source
detectable at feasible array sizes ($N \leq 10^6$) is IMBH binary
mergers. The P6 array sensitivity is $h_{\rm min} \sim 3\ee{-16}
/\sqrt{\rm Hz}$ at $N = 10^6$ (best case), which is $10^4\times$
worse than LISA at LISA's optimal frequency. P6 value comes from
incremental scalability and technology synergy with funded programmes.

\textbf{Most valuable claim:} Ground-based decihertz GW monitoring
in the 0.01--1~Hz gap, leveraging MAGIS/AION technology.

\textbf{Dependencies:} P6 depends on P9 (sensor technology), P14
(field equations for signal prediction). P6 is largely independent
of the SRF detection programme (P11).
\end{tcolorbox}

%=============================================================================
\section{Patent 7: Psi-Field Energy Storage}
\label{sec:P7}
%=============================================================================

\begin{tcolorbox}[colback=yellow!5!white, colframe=yellow!60!black,
  title=\textbf{P7: NICHE --- Resource Allocation: 2/5}]
\textbf{Status: NICHE (ultra-fast pulsed buffer only).}
With $\tau_{1/e} = 1$~s and $E_{\rm max} = 46$~MJ (baseline Nb SRF),
P7 is eliminated for any application requiring $>$3~s storage (grid
frequency regulation, UPS, EV, solar shift --- all dead). Genuine
advantages survive for three application classes: (a) compact pulsed
DEW/HPM ($650\times$ energy density advantage over capacitors in the
cavity, decisive for volume-constrained platforms if cryogenics can
be miniaturised to $<$5~m$^3$); (b) IFE drivers at 10~Hz ($>$99.9\%
round-trip efficiency at 0.1~s cycle, native sub-$\mu$s discharge);
(c) future SRF accelerator RF fill (natural technical match, 98\%
efficiency at 20~ms). Railgun is marginal at baseline but potentially
competitive with Nb$_3$Sn upgrade. System-level volume ($\sim$110~m$^3$
including cryogenics) erases the cavity density advantage for
applications where total system size matters.

\textbf{Most valuable claim:} 650~MJ/m$^3$ pulsed energy storage with
sub-$\mu$s discharge --- unique energy density for ultra-fast pulsed
applications.

\textbf{Dependencies:} P7 depends on P11 (SRF cavity technology),
P2 (cavity materials). P7 shares the SRF infrastructure with P4/P12.
\end{tcolorbox}

%=============================================================================
\section{Patent 8: Psi-Field Communications}
\label{sec:P8}
%=============================================================================

\begin{tcolorbox}[colback=red!5!white, colframe=red!60!black,
  title=\textbf{P8: DEAD --- Resource Allocation: 0/5}]
\textbf{Status: DEAD.}
W3i7 delivered the final blow: $Q_\psi$ amplification in an SRF cavity
does NOT reduce transmitter mass (radiated power = input power,
independent of $Q$ --- a universal property of driven resonant systems).
Every conceivable modulation mechanism (mass modulation, EM-to-$\psi$
Process~A, coherent superradiance, nuclear energy density, rotating
mass quadrupole) encounters the fundamental $G_N/c^5$ barrier. Minimum
transmitter mass for 1~kbit/s at 1~AU with $10^6$ receivers:
$\sim 10^{25}$~kg (1.6 Earth masses). For submarine at 500~m with
$10^4$ receivers: $\sim 3\ee{17}$~kg (Ceres-mass). Mass gap from
achievable technology: $10^{11}\times$ (submarine) to $10^{19}\times$
(interplanetary). The theoretical uniqueness of matter-penetrating
communications has no path to practical realisation.

\textbf{Most valuable claim:} Theoretical only --- the principle that
$\psi$-field modulation penetrates arbitrary matter.

\textbf{Dependencies:} P8 depends on P11 (detection) and P6/P9
(receiver technology), but neither can close the mass gap.
\end{tcolorbox}

%=============================================================================
\section{Patent 9: Psi-Field Positioning and Navigation}
\label{sec:P9}
%=============================================================================

\begin{tcolorbox}[colback=green!5!white, colframe=green!40!black,
  title=\textbf{P9: ALIVE --- Resource Allocation: 3/5}]
\textbf{Status: ALIVE (transformative at SQL; broad application base).}
W3i7 clarified the SQL vs.\ Heisenberg distinction: the 8.7~nm claim
requires Heisenberg-limited sensors (current squeezing achieves only
10--20~dB, needing $\sim 100\times$ improvement). At the SQL, array
positioning is $\sim 9~\mu$m (100 sensors, $N_a = 10^6$) and single
sensor is 0.091~mm. Crucially, \emph{all major applications survive
at SQL}: submarine navigation ($<$1~m required), underwater positioning
($<$10~cm), geological survey ($<$1~m), seismic monitoring ($<$1~mm),
autonomous vehicle navigation ($<$10~cm), infrastructure monitoring
($<$1~mm). SQL single-sensor precision of 0.091~mm is already
$10{,}000\times$ better than the best acoustic positioning (LBL at
$\sim$1~cm). Geological survey at SQL detects $\sim$18~m diameter
ore bodies at 23~km depth. The key advantage: unlimited range,
matter-penetrating, GPS-independent.

\textbf{Most valuable claim:} GPS-independent, matter-penetrating
sub-mm positioning at unlimited range using gravitational field
gradient mapping.

\textbf{Dependencies:} P9 depends on P14 (field equations), P11
(confirming field existence). P9 sensor technology feeds P6 (GW
detection). P9 is the broadest commercial application.
\end{tcolorbox}

%=============================================================================
\section{Patent 10: Field Computing and Logic}
\label{sec:P10}
%=============================================================================

\begin{tcolorbox}[colback=yellow!5!white, colframe=yellow!60!black,
  title=\textbf{P10: NICHE --- Resource Allocation: 1/5}]
\textbf{Status: NICHE (contingent on P3 and P11).}
BQP-completeness is confirmed: psi-field computing can simulate any
quantum circuit. Practical speedup factors are standard BQP advantages
(not DFD-specific): $\sim 10^{12}\times$ for quantum simulation,
$\sim 10^{15}\times$ for cryptography (Shor), $\sqrt{N}$ for search
(Grover), and 2--$100\times$ heuristic for optimisation. The DFD
contribution is the \emph{physical resource reduction} from P3
(10--28$\times$ fewer qubits, 12--32$\times$ shallower circuits),
not new computational complexity advantages. All speedups are
contingent on building a psi-field quantum computer, which requires
P3 (QEC codes), P11 (detection/coupling), and P2 (cavity technology).
No independent experimental programme is needed; P10 rides on P3/P11
progress.

\textbf{Most valuable claim:} Physical resource reduction for quantum
computing via psi-QEC gravitational noise immunity.

\textbf{Dependencies:} P10 depends on P3 (QEC), P11 (coupling), P2
(cavities). P10 is downstream of all hardware patents.
\end{tcolorbox}

%=============================================================================
\section{Patent 11: EM Coupling / SRF Cavity Detection}
\label{sec:P11}
%=============================================================================

\begin{tcolorbox}[colback=green!5!white, colframe=green!40!black,
  title=\textbf{P11: ALIVE --- Resource Allocation: 5/5 (HIGHEST PRIORITY)}]
\textbf{Status: ALIVE (gateway experiment for all of DFD).}
P11 is the most critical patent in the portfolio. The three-phase
detection roadmap is fully specified:
Phase~I: SQMS 4-cavity incoherent array at Fermilab (\$2--4M,
2027--2028, sensitivity $\lambdaone = 7.8\ee{-10}$);
Phase~II: MEMS hybrid with Si$_3$N$_4$ phononic crystal membrane
(2028--2030, $\lambdaone = 4.7\ee{-11}$);
Phase~III: 256-cavity coherent array or ng-MEMS (2030--2035,
$\lambdaone \sim 10^{-12}$).
W3i7 provided the complete SQMS Phase~I protocol, systematic error
budget, and the indirect $Q_\psi$ constraint from existing multi-frequency
data. The current bound ($\lambdaone < 1.56\ee{-9}$) comes from SQMS
passive $Q_0$ measurements. Phase~I improves this by $2\times$; Phase~II
by $30\times$; Phase~III by $>1000\times$. Detection at any phase
activates the entire downstream patent portfolio. Null result at
Phase~III would require fundamental revision.

\textbf{Most valuable claim:} SRF cavity as a $\psi$-field detector
via Process~A ($\lambdaone$-dependent $Q_0$ anomaly).

\textbf{Dependencies:} P11 depends on P2 (Nb cavity technology), P14
(coupling physics). P11 $\to$ ALL other patents (gateway).
\end{tcolorbox}

%=============================================================================
\section{Patent 12: City-Scale Psi-Transmission Network}
\label{sec:P12}
%=============================================================================

\begin{tcolorbox}[colback=orange!5!white, colframe=orange!60!black,
  title=\textbf{P12: THEORETICAL ONLY --- Resource Allocation: 0/5}]
\textbf{Status: THEORETICAL ONLY.}
The city-scale psi-transmission network is economically non-viable at
20\% end-to-end efficiency. The LCOD is dominated by energy losses:
$\text{LCOD} \approx 5 \times p_{\rm elec}$, giving 25~c/kWh at
\$0.05/kWh input. This is $2.5\times$ US retail average. Capital cost
is irrelevant ($<$1\% of LCOD at city-scale throughput). W3i7 identified
the critical correction: P4 and P12 share the same conversion bottleneck
(two 50\% stages). The narrow viability window requires ultra-cheap
input ($<$\$0.02/kWh) AND extremely expensive conventional delivery
($>$\$0.25/kWh). Niche markets (deep undersea, deep mining, space)
are theoretically viable but are better served by P4 point-to-point
corridors. No resources should be allocated to P12 until the
conversion bottleneck is broken ($\eta_{\rm e2e} > 50\%$).

\textbf{Most valuable claim:} Network topology analysis for distributed
psi-power delivery (deferred value).

\textbf{Dependencies:} P12 depends on P4 (corridor technology), P11
(detection), P15/P2 (conversion efficiency). P12 activates only if
conversion efficiency exceeds 50\%.
\end{tcolorbox}

%=============================================================================
\section{Patent 13: Clock Comparison Tests}
\label{sec:P13}
%=============================================================================

\begin{tcolorbox}[colback=yellow!5!white, colframe=yellow!60!black,
  title=\textbf{P13: NICHE --- Resource Allocation: 3/5}]
\textbf{Status: NICHE (high near-term value for theory validation).}
P13 provides the most accessible near-term experimental test of DFD
via clock comparison measurements. CLONETS-DS achieves detection in
1.0~day and sidereal discrimination at 87-sigma in 430~days. W3i7
delivered the complete measurement protocol with five phases
(commissioning, detection, characterisation, sidereal confirmation,
precision programme), full systematic error budget, and null tests.
Five archival predictions are testable \emph{now} from existing
public datasets: GPS diurnal drift, Cassini ranging residuals,
pulsar timing, galaxy rotation curve shapes, and Tully--Fisher
zero-point. These archival tests are essentially free and could
provide early evidence before any new experiment. P13 synergises
with P5 (tests the same underlying predictions) and is independent
of the SRF programme.

\textbf{Most valuable claim:} Sub-day detection of DFD clock effects
using existing optical clock network infrastructure.

\textbf{Dependencies:} P13 depends on P14 (prediction framework), P5
(prediction catalogue). P13 is independent of P11 (different detection
method). P13 $\to$ P5 (validates cosmological sector).
\end{tcolorbox}

%=============================================================================
\section{Patent 14: Canonical Two-Component Framework (ICC)}
\label{sec:P14}
%=============================================================================

\begin{tcolorbox}[colback=green!5!white, colframe=green!40!black,
  title=\textbf{P14: ALIVE --- Resource Allocation: 4/5}]
\textbf{Status: ALIVE (foundational; underpins everything).}
P14 provides the canonical formulation: 4 axioms, 2 free parameters
($\lambda$, $\kappa$), 13+ testable predictions. W3i7 confirmed that
$\alpha^{57}$ (the vacuum energy chain) is unaffected by the
two-component decomposition (UV/IR separation argument); the
$\mu(x) = x/(1+x)$ MOND function is derived from the $S^3$ composition
law and is compatible with the two-component framework; $\kappa$ is
operationally irrelevant at current sensitivity, effectively leaving
one free parameter ($\lambda$). The T4 Cold Spot over-prediction was
reduced from $3$--$5\times$ to $1$--$2\times$ by the void-evolution
cancellation. The MOND scale $a_0 = 2\sqrt{\alpha}\,cH_0$ is a
zero-parameter prediction (corrected from the W3i6 $\sqrt{\Omega_\Lambda}$
formula). P14 is not an independent ``product'' but the theoretical
engine driving all other patents. Resource allocation here means
theoretical development and publication.

\textbf{Most valuable claim:} Single scalar field unifying gravity,
MOND, dark energy, and EM coupling with one effective free parameter.

\textbf{Dependencies:} P14 depends on nothing (foundational). P14
$\to$ ALL other patents.
\end{tcolorbox}

%=============================================================================
\section{Patent 15: Psi-Field Generator / Detection Channel B}
\label{sec:P15}
%=============================================================================

\begin{tcolorbox}[colback=yellow!5!white, colframe=yellow!60!black,
  title=\textbf{P15: NICHE --- Resource Allocation: 3/5}]
\textbf{Status: NICHE (highest experimental priority after P11).}
P15 Channel~B (passive $Q$-anomaly measurement at ESS Lund) is the
highest-priority experimental channel identified in W3i7. Channel~A
(direct $\dpsiEM$ detection) cannot improve the current $\lambdaone$
bound but remains valuable for systematics verification. The ESS
measurement protocol is fully specified: three-phase protocol with
explicit 5-sigma criteria including frequency discrimination,
temperature independence, and universality tests. Cost is modest
(\$50K--\$100K for beam time and instrumentation). P15 Channel~B
requires P2 toroidal cavity technology and P11 SRF coating, creating
a strong synergy. The P15 generator concept (actively producing
$\dpsiEM$) is blocked by the same $G_N/c^4$ coupling weakness that
killed P8, but Channel~B (passive detection in existing particle
beams) circumvents this by using existing high-energy sources.

\textbf{Most valuable claim:} Passive $Q$-anomaly detection in SRF
cavities exposed to particle beams (ESS, SNS, J-PARC).

\textbf{Dependencies:} P15 depends on P2 (cavity technology), P11
(SRF expertise). P15 $\to$ P11 (alternative detection pathway).
\end{tcolorbox}


%=============================================================================
%=============================================================================
\part{Summary Tables}
\label{part:tables}
%=============================================================================
%=============================================================================

%=============================================================================
\section{Master Patent Status Table}
\label{sec:master-table}
%=============================================================================

\begin{landscape}
\begin{longtable}{@{}clp{4cm}p{3.5cm}cp{3.5cm}@{}}
\caption{W3i8 Definitive Patent Triage --- All 15 DFD Patents.
Resource allocation on 0--5 scale.}
\label{tab:master}\\
\toprule
\textbf{Patent} & \textbf{Status} & \textbf{What Survives}
  & \textbf{Most Valuable Claim} & \textbf{Resources}
  & \textbf{Key Dependencies} \\
\midrule
\endfirsthead
\multicolumn{6}{c}{\textit{Table~\ref{tab:master} continued}}\\
\toprule
\textbf{Patent} & \textbf{Status} & \textbf{What Survives}
  & \textbf{Most Valuable Claim} & \textbf{Resources}
  & \textbf{Key Dependencies} \\
\midrule
\endhead
\bottomrule
\endlastentry
P1 & \textcolor{red}{DEAD} & Metric structure (to P14) &
  DFD metric coupling & 0 & None active \\[4pt]
P2 & \textcolor{orange}{NICHE} & Passive detection; P15 cavity &
  Nb SRF as $\psi$-detector & 2 & P11, P14 \\[4pt]
P3 & \textcolor{orange}{NICHE} & 7.7--83$\times$ qubit advantage &
  $67^K$ error suppression & 2 & P11 (gate) \\[4pt]
P4 & \textcolor{orange}{NICHE} & Matter-penetrating corridor &
  Underground WPT & 2 & P11, conversion $\eta$ \\[4pt]
P5 & \textcolor{green!50!black}{ALIVE} & 11/13 predictions; Hubble
  tension & Unified cosmology & 4 & P14 \\[4pt]
P6 & \textcolor{orange}{NICHE} & Ground-based decihertz access &
  0.01--1 Hz gap-filler & 1 & P9 sensors \\[4pt]
P7 & \textcolor{orange}{NICHE} & DEW/fusion/accelerator buffer &
  650 MJ/m$^3$ pulsed storage & 2 & P11 SRF \\[4pt]
P8 & \textcolor{red}{DEAD} & Theoretical principle only &
  Matter-penetrating comms & 0 & Unfeasible \\[4pt]
P9 & \textcolor{green!50!black}{ALIVE} & SQL sub-mm positioning &
  GPS-free navigation & 3 & P14, P11 \\[4pt]
P10 & \textcolor{orange}{NICHE} & BQP resource reduction &
  Physical qubit savings & 1 & P3, P11 \\[4pt]
P11 & \textcolor{green!50!black}{ALIVE} & Gateway experiment &
  SRF $\psi$-detector & \textbf{5} & P2, P14 \\[4pt]
P12 & \textcolor{orange!70!black}{THEORY} & Network topology analysis &
  Distributed psi-power & 0 & P4, conversion $\eta$ \\[4pt]
P13 & \textcolor{orange}{NICHE} & CLONETS-DS; archival tests &
  1-day clock detection & 3 & P14, P5 \\[4pt]
P14 & \textcolor{green!50!black}{ALIVE} & Canonical framework &
  Unified scalar field & 4 & None (foundation) \\[4pt]
P15 & \textcolor{orange}{NICHE} & Channel B at ESS &
  Passive $Q$-anomaly & 3 & P2, P11 \\
\bottomrule
\end{longtable}
\end{landscape}

%=============================================================================
\section{Resource Allocation Summary}
\label{sec:resources}
%=============================================================================

\begin{table}[H]
\centering
\caption{Resource allocation by priority tier. Total budget units: 32/75
(57\% of maximum, reflecting 2 dead + 1 theoretical-only patents).}
\label{tab:resources}
\begin{tabular}{@{}clcc@{}}
\toprule
\textbf{Tier} & \textbf{Patents} & \textbf{Allocation} &
  \textbf{Total Units} \\
\midrule
\multirow{4}{*}{1 --- ALIVE} & P11 (Detection) & 5/5 & \\
  & P5 (Cosmology) & 4/5 & \\
  & P14 (Framework) & 4/5 & \\
  & P9 (Positioning) & 3/5 & 16 \\
\midrule
\multirow{8}{*}{2 --- NICHE} & P13 (Clocks) & 3/5 & \\
  & P15 (Generator/B) & 3/5 & \\
  & P2 (Resonators) & 2/5 & \\
  & P3 (QEC) & 2/5 & \\
  & P4 (Underground WPT) & 2/5 & \\
  & P7 (Storage) & 2/5 & \\
  & P10 (Computing) & 1/5 & \\
  & P6 (Decihertz GW) & 1/5 & 16 \\
\midrule
3 --- THEORY & P12 (City WPT) & 0/5 & 0 \\
\midrule
\multirow{2}{*}{4 --- DEAD} & P1 (Field Drag) & 0/5 & \\
  & P8 (Comms) & 0/5 & 0 \\
\bottomrule
\end{tabular}
\end{table}

%=============================================================================
\section{Dependency Graph}
\label{sec:dependency}
%=============================================================================

The patent dependency structure is hierarchical with P14 at the foundation
and P11 as the experimental gateway:

\begin{verbatim}
                        P14 (Framework)
                       /    |    \    \
                      /     |     \    \
                    P5     P11    P9   P13
                    |     / | \    |
                   P13  P2 P15 P7  P6
                        |       |
                       P4      P3
                        |       |
                       P12    P10
\end{verbatim}

\begin{finding}[Critical Path]
\label{find:critical-path}
The critical path through the DFD programme is:
\begin{equation*}
\text{P14} \to \text{P11 Phase~I} \to
\begin{cases}
\text{P11 Phase~II} \to \text{P15 Channel~B} \to \text{P4/P7 prototypes}\\
\text{P5 archival validation} \to \text{P13 CLONETS-DS}\\
\text{P9 sensor development} \to \text{P6 array deployment}
\end{cases}
\end{equation*}

\textbf{Single point of failure}: P11. If P11 Phase~III returns null,
the entire technology portfolio (P1--P4, P7--P8, P10, P12) is invalidated.
Only P5 (cosmological predictions), P9 (gravity gradient positioning,
which works regardless of $\dpsiEM$), P13 (clock tests), and P14
(theoretical framework) survive a null P11 result.
\end{finding}


%=============================================================================
%=============================================================================
\part{Complete Technology Roadmap}
\label{part:roadmap}
%=============================================================================
%=============================================================================

%=============================================================================
\section{Phase 0: Archival Analysis (2026, \$0, immediate)}
\label{sec:phase0}
%=============================================================================

\begin{tcolorbox}[colback=green!5!white, colframe=green!40!black,
  title=\textbf{Phase 0: Zero-Cost Archival Tests}]

\textbf{What gets built:} Nothing. Pure data analysis.

\textbf{Cost:} \$0 (researcher time only).

\textbf{Timeline:} 2026, starting immediately.

\textbf{What it tests:} Five P5 predictions from existing public datasets.

\begin{enumerate}[nosep]
\item \textbf{GPS diurnal drift (P5/P13):} Reanalyse IGS precise clock
  products for diurnal $\psigrav$ modulation at $\sim 10^{-16}$ fractional
  frequency. Data: freely available from IGS archives (20+ years).
  Expected signal: diurnal sinusoid with amplitude set by
  $\Delta\psigrav = GM_\oplus/(c^2 R_\oplus) \approx 7\ee{-10}$
  modulated by Earth's rotation relative to galactic centre.

\item \textbf{Cassini ranging residuals (P5):} Reanalyse Cassini
  radiometric data for anomalous frequency drift consistent with
  cosmological $\psi$-evolution. Data: NASA PDS. Expected signal:
  $\dot{\nu}/\nu \sim H_0 \sim 10^{-18}$~s$^{-1}$ (at the noise floor
  but constrainable with 13~years of data).

\item \textbf{Pulsar timing array residuals (P5):} Search IPTA data
  release 2 for monopole/dipole correlations consistent with scalar
  breathing mode (P2 LIGO 6th channel, but at nHz frequencies).
  Data: IPTA DR2, freely available.

\item \textbf{Galaxy rotation curve shape parameter (P5):} Test the
  DFD prediction that rotation curves follow $\mu(g/a_0)$ with
  $a_0 = 2\sqrt{\alpha}\,cH_0 \approx 1.1\ee{-10}$~m/s$^2$
  using SPARC database (175 galaxies, public).

\item \textbf{Tully--Fisher zero-point (P5):} Test the predicted
  baryonic Tully--Fisher relation slope and zero-point against the
  McGaugh (2012) dataset and more recent compilations.
\end{enumerate}

\textbf{Success criteria:} Any single archival detection at $>3\sigma$
justifies accelerating Phase~I funding.
\end{tcolorbox}

%=============================================================================
\section{Phase I: First Detection Experiments (2027--2028, \$2--5M)}
\label{sec:phase1}
%=============================================================================

\begin{tcolorbox}[colback=blue!5!white, colframe=blue!60!black,
  title=\textbf{Phase I: SQMS 4-Cavity + ESS Channel B}]

\textbf{What gets built:}

\begin{enumerate}
\item \textbf{SQMS 4-cavity incoherent array (P11 Phase~I):}
  \begin{itemize}[nosep]
    \item Location: Fermilab SQMS Centre (existing infrastructure).
    \item Hardware: 4 TESLA-type Nb SRF cavities, medium-$T$ N-doped,
      5~$\mu$m EP, 120$^\circ$C/48h bake. Sub-nT magnetic shielding.
    \item Operating temperature: 2~K (He-II).
    \item Measurement: Multi-frequency $Q_0(T)$ curves, searching for
      $\lambdaone$-dependent anomaly in the temperature-dependent
      dissipation.
    \item Sensitivity: $\lambdaone = 7.8\ee{-10}$ (incoherent average
      of 4 cavities, $2\times$ improvement over current bound).
    \item Cost: \$2--4M (cavities, cryogenics, shielding, 2~years FTE).
    \item Timeline: 2027--2028.
  \end{itemize}

\item \textbf{ESS Lund passive Q-anomaly (P15 Channel~B):}
  \begin{itemize}[nosep]
    \item Location: ESS Lund superconducting linac (84 SRF cavities
      in operation).
    \item Measurement: $Q_0$ monitoring during beam-on/beam-off cycles.
      Three-phase protocol: frequency discrimination (comparing
      cavities at different frequencies), temperature independence
      check, universality test across cavity types.
    \item Sensitivity: $\lambdaone_{\rm min} \sim 10^{-7}$ (single
      cavity); $\sim 10^{-7}$ (statistical, 84 cavities). Cannot
      improve the current bound but tests different systematics.
    \item Cost: \$50K--\$100K (beam time, instrumentation).
    \item Timeline: 2027 (parasitic on ESS commissioning).
  \end{itemize}
\end{enumerate}

\textbf{What it tests:}
\begin{itemize}[nosep]
  \item Does $\lambdaone > 0$ at $7.8\ee{-10}$ sensitivity? (P11)
  \item Is the $Q_0$ anomaly universal across cavity frequencies and
    temperatures? (P15)
  \item Can existing SQMS data already constrain $Q_\psi$? (P11 indirect)
\end{itemize}

\textbf{Decision gate:}
\begin{itemize}[nosep]
  \item \textbf{Detection ($>5\sigma$):} Proceed to Phase~II at maximum
    speed. Activate P4/P7 prototype planning.
  \item \textbf{Null ($\lambdaone < 7.8\ee{-10}$):} Proceed to Phase~II
    (MEMS hybrid) for deeper sensitivity. No technology patents activated.
  \item \textbf{Anomalous systematic:} Investigate. Redesign Phase~II
    to address.
\end{itemize}
\end{tcolorbox}

%=============================================================================
\section{Phase II: Deep Detection and Sensor Development (2028--2031, \$5--20M)}
\label{sec:phase2}
%=============================================================================

\begin{tcolorbox}[colback=blue!5!white, colframe=blue!60!black,
  title=\textbf{Phase II: MEMS Hybrid + P9 Sensor Prototype}]

\textbf{What gets built:}

\begin{enumerate}
\item \textbf{MEMS + SRF hybrid detector (P11 Phase~II / P2):}
  \begin{itemize}[nosep]
    \item Si$_3$N$_4$ soft-clamped phononic crystal membrane:
      1~mm $\times$ 1~mm $\times$ 20~nm, $\sigma = 1.2$~GPa,
      $f_0 \approx 1.4$~MHz, $m_{\rm eff} = 0.6$~ng,
      $Q_{\rm mech} \geq 10^9$ (room temperature, demonstrated;
      mK projected $>10^{10}$).
    \item Coupled to SRF cavity readout.
    \item Sensitivity: $\lambdaone = 4.7\ee{-11}$ ($30\times$
      improvement over Phase~I).
    \item Cost: \$5--10M.
    \item Timeline: 2028--2030.
  \end{itemize}

\item \textbf{P9 positioning sensor prototype:}
  \begin{itemize}[nosep]
    \item Single atom interferometer sensor, MAGIS/AION-derived.
    \item Target: demonstrate 0.091~mm single-sensor positioning
      (SQL) in laboratory setting.
    \item Test underwater operation (confirm $\psi$-field immunity
      to water column).
    \item Cost: \$3--5M (leveraging MAGIS/AION technology).
    \item Timeline: 2029--2031.
  \end{itemize}

\item \textbf{CLONETS-DS pathfinder (P13):}
  \begin{itemize}[nosep]
    \item Deploy optical clock comparison on existing fibre links
      (European optical clock network).
    \item 1-day detection sensitivity; 430-day sidereal confirmation.
    \item Cost: \$1--3M (incremental to existing network).
    \item Timeline: 2028--2030.
  \end{itemize}
\end{enumerate}

\textbf{What it tests:}
\begin{itemize}[nosep]
  \item Does $\lambdaone > 0$ at $4.7\ee{-11}$? (P11)
  \item Can gravitational gradient mapping achieve sub-mm positioning?
    (P9)
  \item Is there a sidereal modulation in clock comparisons? (P13/P5)
\end{itemize}

\textbf{Decision gate:}
\begin{itemize}[nosep]
  \item \textbf{P11 detection:} Activate P4 corridor prototype, P7
    storage prototype, P3/P10 quantum computing planning.
  \item \textbf{P11 null + P13 detection:} DFD validated via
    gravitational sector only; $\dpsiEM$ coupling weaker than expected.
    Revise $\lambda$ bounds. Continue to Phase~III.
  \item \textbf{Both null:} Proceed to Phase~III but with reduced
    confidence. Begin planning for ``DFD falsification'' scenario.
\end{itemize}
\end{tcolorbox}

%=============================================================================
\section{Phase III: Technology Prototypes (2031--2035, \$50--200M)}
\label{sec:phase3}
%=============================================================================

\textbf{Contingent on Phase~II detection in at least one channel.}

\begin{tcolorbox}[colback=blue!5!white, colframe=blue!60!black,
  title=\textbf{Phase III: First Technology Demonstrations}]

\begin{enumerate}
\item \textbf{256-cavity coherent array (P11 Phase~III):}
  \begin{itemize}[nosep]
    \item Sensitivity: $\lambdaone \sim 10^{-12}$ ($>1000\times$
      current bound).
    \item Facility-class: 2--5~MW cryogenic power.
    \item Cost: \$20--50M.
    \item Tests: Precision measurement of $\lambda$; characterisation
      of $\dpsiEM$ spectrum; direct $Q_\psi$ measurement.
  \end{itemize}

\item \textbf{P4 underground WPT corridor prototype:}
  \begin{itemize}[nosep]
    \item 1~km underground corridor (mine or tunnel).
    \item SRF transmitter + SRF receiver, each at 50\% conversion.
    \item Target: demonstrate end-to-end wireless power transfer
      through rock.
    \item Cost: \$20--50M.
    \item Tests: Transport efficiency vs.\ distance; conversion
      efficiency at SRF; $Q_\psi$ along corridor.
  \end{itemize}

\item \textbf{P7 pulsed storage demonstrator:}
  \begin{itemize}[nosep]
    \item Single Nb$_3$Sn SRF cavity, 4~K operation.
    \item Target: 10~MJ stored, $<$1~ms discharge, $>$99\% round-trip.
    \item Cost: \$5--15M.
    \item Tests: Energy density, discharge rate, $\tau_{1/e}$
      measurement.
  \end{itemize}

\item \textbf{P9 multi-sensor array:}
  \begin{itemize}[nosep]
    \item 10--100 sensor array, fixed installation.
    \item Target: demonstrate $\sim 10~\mu$m array positioning (SQL).
    \item Underwater deployment for submarine navigation proof of
      concept.
    \item Cost: \$10--30M.
    \item Tests: Array scaling ($1/\sqrt{N}$); underwater performance;
      geological mapping resolution.
  \end{itemize}

\item \textbf{P6 decihertz pathfinder:}
  \begin{itemize}[nosep]
    \item Piggyback on MAGIS-100 or AION-10.
    \item 10-sensor array at existing facility.
    \item Cost: \$2--5M (incremental).
    \item Tests: Decihertz noise characterisation; $\psi$-field
      sensitivity at 0.01--1~Hz.
  \end{itemize}
\end{enumerate}
\end{tcolorbox}

%=============================================================================
\section{Phase IV: Scaling and Commercialisation (2035+, \$200M--\$1B+)}
\label{sec:phase4}
%=============================================================================

\textbf{Contingent on Phase~III successful demonstrations.}

\begin{tcolorbox}[colback=blue!5!white, colframe=blue!60!black,
  title=\textbf{Phase IV: Full-Scale Deployment}]

\begin{enumerate}
\item \textbf{P9 commercial positioning system:}
  \begin{itemize}[nosep]
    \item First commercial product: underwater/underground GPS-free
      navigation.
    \item Target markets: submarine navigation, mining, offshore oil,
      autonomous vehicles.
    \item Revenue potential: \$1B+ (GPS-denied navigation market).
  \end{itemize}

\item \textbf{P4 undersea power corridor:}
  \begin{itemize}[nosep]
    \item First niche deployment: deep-sea power delivery to ocean
      floor installations.
    \item Viable when conventional undersea cable costs $>$\$5M/km.
    \item Cost: \$80--200M per installation.
  \end{itemize}

\item \textbf{P7 military/fusion pulsed storage:}
  \begin{itemize}[nosep]
    \item Compact DEW power supply (platform-optimised cryogenics).
    \item IFE driver for fusion energy plants.
  \end{itemize}

\item \textbf{P3/P10 quantum computing:}
  \begin{itemize}[nosep]
    \item Psi-QEC quantum processor (requires mature SRF + MEMS
      technology from P11/P2).
    \item Advantage: 10--28$\times$ fewer qubits than surface codes.
  \end{itemize}

\item \textbf{P6 decihertz GW observatory:}
  \begin{itemize}[nosep]
    \item $10^4$--$10^6$ sensor array for IMBH binary detection.
    \item Timeline: 2040+.
  \end{itemize}
\end{enumerate}
\end{tcolorbox}


%=============================================================================
%=============================================================================
\part{Roadmap Summary and Critical Milestones}
\label{part:milestones}
%=============================================================================
%=============================================================================

%=============================================================================
\section{Timeline Overview}
\label{sec:timeline}
%=============================================================================

\begin{table}[H]
\centering
\caption{Complete DFD technology roadmap timeline.}
\label{tab:timeline}
\small
\begin{tabular}{@{}clp{5cm}cp{3cm}@{}}
\toprule
\textbf{Phase} & \textbf{Year} & \textbf{Experiment} & \textbf{Cost} &
  \textbf{Patents Tested} \\
\midrule
\multirow{5}{*}{0} & \multirow{5}{*}{2026} & GPS diurnal reanalysis &
  \$0 & P5, P13 \\
  & & Cassini residuals & \$0 & P5 \\
  & & Pulsar timing search & \$0 & P5, P2 \\
  & & SPARC rotation curves & \$0 & P5, P14 \\
  & & Tully--Fisher test & \$0 & P5 \\
\midrule
\multirow{2}{*}{I} & 2027--28 & SQMS 4-cavity (P11) & \$2--4M &
  P11, P2 \\
  & 2027 & ESS Channel~B (P15) & \$50--100K & P15, P11 \\
\midrule
\multirow{3}{*}{II} & 2028--30 & MEMS hybrid (P11/P2) & \$5--10M &
  P11, P2, P3 \\
  & 2029--31 & P9 sensor prototype & \$3--5M & P9 \\
  & 2028--30 & CLONETS-DS (P13) & \$1--3M & P13, P5 \\
\midrule
\multirow{5}{*}{III} & 2031--35 & 256-cavity array & \$20--50M &
  P11 \\
  & & P4 corridor prototype & \$20--50M & P4 \\
  & & P7 storage demo & \$5--15M & P7 \\
  & & P9 multi-sensor array & \$10--30M & P9 \\
  & & P6 decihertz pathfinder & \$2--5M & P6 \\
\midrule
IV & 2035+ & Commercial deployment & \$200M--1B+ &
  P9, P4, P7, P3/P10 \\
\bottomrule
\end{tabular}
\end{table}

%=============================================================================
\section{Critical Milestones and Decision Gates}
\label{sec:milestones}
%=============================================================================

\begin{finding}[Six Make-or-Break Milestones]
\label{find:milestones}
\begin{enumerate}
\item \textbf{M1: Archival evidence (2026).} Any $>3\sigma$ signal in
  GPS diurnal, rotation curves, or Tully--Fisher data validates the P5
  cosmological sector and justifies Phase~I acceleration. Cost: \$0.

\item \textbf{M2: SQMS Phase~I detection (2028).} $\lambdaone$ signal
  at $7.8\ee{-10}$ sensitivity. This is the \textbf{single most important
  milestone}. Detection activates the entire technology portfolio. Null
  tightens the bound by $2\times$ and motivates Phase~II.

\item \textbf{M3: MEMS hybrid detection (2030).} $\lambdaone$ signal at
  $4.7\ee{-11}$. Detection at this depth proves the coupling is real and
  measurable. Null at Phase~II + null at CLONETS-DS would be a serious
  blow to DFD.

\item \textbf{M4: Conversion efficiency breakthrough ($>$50\%).}
  Breaking the two-stage 50\% SRF conversion bottleneck. This single
  milestone transforms P4 from niche (\$0.27/kWh) to competitive
  (\$0.10/kWh) and activates P12. Requires either: (a) direct
  SRF-to-SRF psi-mode transfer, (b) higher $\lambdaone$ (increasing
  Process~A efficiency), or (c) a new conversion mechanism. Timeline
  is unpredictable; this is the hardest open problem.

\item \textbf{M5: P9 sub-mm positioning demonstrated (2031).}
  Laboratory demonstration of 0.091~mm single-sensor gravitational
  gradient positioning. This is the first commercially relevant
  milestone and the gateway to the GPS-denied navigation market.

\item \textbf{M6: P4 underground power transfer demonstrated (2035).}
  End-to-end wireless power through rock. Even at 20\% efficiency,
  this demonstrates a capability that no other physics can provide
  (wireless power through arbitrary matter), opening niche markets.
\end{enumerate}
\end{finding}

%=============================================================================
\section{Cumulative Investment Profile}
\label{sec:investment}
%=============================================================================

\begin{table}[H]
\centering
\caption{Cumulative investment and expected return milestones.}
\label{tab:investment}
\begin{tabular}{@{}ccll@{}}
\toprule
\textbf{Year} & \textbf{Cumulative \$} & \textbf{What is known} &
  \textbf{What is activated} \\
\midrule
2026 & \$0 & Archival evidence (if any) & Phase~I justification \\
2028 & \$2--5M & $\lambdaone$ bound at $7.8\ee{-10}$ &
  Phase~II or technology activation \\
2031 & \$10--25M & Detection or null at $4.7\ee{-11}$; P9 demo &
  Phase~III or reassessment \\
2035 & \$70--175M & Full characterisation; prototypes &
  Commercial planning \\
2040+ & \$200M--1B & Commercial products & Revenue \\
\bottomrule
\end{tabular}
\end{table}

%=============================================================================
\section{The Conversion Bottleneck: The Single Most Important Technology Problem}
\label{sec:bottleneck}
%=============================================================================

\begin{tcolorbox}[colback=red!5!white, colframe=red!60!black,
  title=\textbf{THE BINDING CONSTRAINT: 50\% SRF Conversion Efficiency}]

The SRF polariton-crossing conversion achieves $\sim$50\% per stage.
Two stages (EM $\to$ $\dpsiEM$ $\to$ EM) give 25\% conversion
efficiency. With 81\% transport, the end-to-end is $\sim$20\%.

\textbf{This single bottleneck determines the economics of P4, P7,
and P12.}

\begin{center}
\begin{tabular}{@{}ccl@{}}
\toprule
$\eta_{\rm conv}$ (per stage) & $\eta_{\rm e2e}$ & Economic status \\
\midrule
50\% (current) & 20\% & Niche only ($>$25~c/kWh) \\
70\% & 40\% & Competitive in some markets ($\sim$13~c/kWh) \\
80\% & 52\% & Broadly competitive ($\sim$10~c/kWh) \\
90\% & 66\% & Grid-competitive ($\sim$8~c/kWh) \\
Direct psi-mode & 81\% & Cheaper than conventional ($\sim$6~c/kWh) \\
\bottomrule
\end{tabular}
\end{center}

\textbf{Paths to breaking the bottleneck:}
\begin{enumerate}[nosep]
\item \textbf{Higher $\lambdaone$:} If $\lambdaone$ is larger than
  the current bound, Process~A conversion is inherently more efficient.
  P11 Phase~I will constrain this.
\item \textbf{Direct SRF-to-SRF transfer:} Bypass the polariton
  crossing; transfer $\dpsiEM$ directly between resonant cavities.
  Eliminates one conversion stage.
\item \textbf{Nb$_3$Sn cavities:} Higher $B_{\rm sh}$ (425~mT vs.\
  180~mT for Nb) increases $P_{\rm max}$ by $\sim 5.6\times$, but
  does not directly improve conversion efficiency.
\item \textbf{Psi-mode generator:} Generate $\dpsiEM$ directly from
  an energy source without EM intermediary. Eliminates both conversion
  stages. Physics of this is an open question.
\end{enumerate}
\end{tcolorbox}

%=============================================================================
\section{What Survives a Null P11 Result?}
\label{sec:null-scenario}
%=============================================================================

\begin{finding}[Null Scenario: DFD Without $\dpsiEM$ Coupling]
\label{find:null}
If P11 Phase~III returns null ($\lambdaone < 10^{-12}$), the following
patents survive:

\begin{itemize}[nosep]
\item \textbf{P5 (Cosmology):} All gravitational-sector predictions
  ($\psigrav$-based) survive. Hubble tension, MOND, dark energy, galaxy
  rotation curves are all $\psigrav$ effects. P5 is the most robust
  patent against a null P11 result.
\item \textbf{P9 (Positioning):} Gravitational gradient mapping uses
  $\psigrav$, not $\dpsiEM$. GPS-free navigation works regardless.
  However, the $10{,}000\times$ advantage over acoustic positioning
  comes from the DFD-predicted noise immunity, which \emph{does} depend
  on the psi-field structure. At minimum, standard gravity gradiometry
  (existing technology) provides $\sim$1~mGal sensitivity.
\item \textbf{P13 (Clocks):} Clock comparison tests measure $\psigrav$
  effects. Survive regardless of $\dpsiEM$.
\item \textbf{P14 (Framework):} The gravitational sector ($\psigrav$
  equation) remains valid even if $\dpsiEM$ coupling is undetectably
  small. DFD reduces to a modified gravity theory.
\end{itemize}

\textbf{Dead in null scenario:} P1, P2, P3, P4, P7, P8, P10, P11,
P12, P15 --- all technology patents depending on $\dpsiEM$ coupling.

\textbf{Assessment:} Even in the worst case, DFD retains scientific
value as a modified gravity framework with cosmological predictions.
The technology portfolio is lost, but the theoretical framework may
still be correct with $\lambdaone$ below any foreseeable detection
threshold.
\end{finding}


%=============================================================================
\section*{Concluding Assessment}
%=============================================================================

After seven iterations of adversarial analysis by 17 parallel agents,
the DFD patent portfolio crystallises into a clear picture:

\textbf{The theory is robust.} P14's canonical framework, P5's
cosmological predictions, and P14's $\alpha^{57}$ chain have survived
every attack. The two-component decomposition resolved rather than
created problems.

\textbf{The technology is gated by one experiment.} P11 SQMS Phase~I
(\$2--4M, 2027--2028) is the single point that determines whether the
technology portfolio has any future. This is not optional expenditure;
it is the minimum viable test of the entire programme.

\textbf{The economics are gated by one bottleneck.} The 50\% SRF
conversion efficiency kills the economics of every energy patent (P4,
P7, P12). Breaking this bottleneck is the hardest open problem and
the highest-leverage technology target.

\textbf{The near-term commercial opportunity is P9.} GPS-denied
sub-mm positioning requires only the gravitational sector, works at
SQL, and addresses a multi-billion-dollar market (underwater, mining,
autonomous vehicles). P9 should be pursued in parallel with P11,
as it is partially independent of the $\dpsiEM$ detection question.

\textbf{The immediate action is Phase~0.} Five archival tests at zero
cost, starting now. If any produces a $>3\sigma$ result, the entire
programme is transformed.


\end{document}
